\section{Summary and Conclusion}\label{summary-and-conclusion}

\subsection{Overview}

This study designed and evaluated GBBRPM v1 as a bounded recursive model
for comparative node-risk prioritization in directed acyclic networks.
The evidence combines controlled architecture tests with three
domain-specific cases: drainage, software dependency, and electrical
networks. The cases intentionally support different claims and are not
presented as equivalent forms of outcome validation.

\subsection{Summary of Findings}

\subsubsection{Architecture Behavior}

Across the controlled N1--N5 experiments, GBBRPM preserved the declared
$[0,1]$ bounds and monotone response to increases in local disturbance
and susceptibility. Recursive current-state propagation, heterogeneous
edge conditions, disturbance location, and multi-source configuration
materially affected both risk magnitude and ranking. Randomized property
tests passed all 600 trials for boundedness, monotonicity, full local
disturbance, and zero-susceptibility gating.

Ranking robustness remained high under modest perturbations, although
agreement among the highest-priority nodes declined faster than the
overall Spearman ordering. The reconvergence diagnostic also showed that
shared ancestry can be counted through multiple paths, especially in N5.
This is a known boundary of the local multiplicative-complement
aggregator rather than evidence of probabilistic independence.

The sparse-DAG timing experiment was consistent with the expected
one-pass $O(|V|+|E|)$ implementation. Median execution time increased
from approximately 1.63 ms at 20 nodes to 16.65 ms at 1000 nodes under
the reported environment.

\subsubsection{Drainage Evidence and Hydraulic Boundary}

The drainage case supplied the most complete behavioral evaluation. The
controlled experiments demonstrated the effects of attenuation,
branching, convergence, reconvergence, local disturbance, susceptibility,
and multiple active sources. However, the 12-scenario SWMM comparison
produced negative mean rank association and low top-20\% overlap for
both maximum-depth and flooding outcomes. Hydraulic worsening also
occurred outside the model's strictly downstream region.

These findings do not support interpreting the GBBRPM index as a
hydraulic predictor. Instead, they identify a clear scope boundary:
GBBRPM represents lightweight downstream comparative propagation,
whereas SWMM represents hydraulic mechanisms such as backwater,
surcharge, ponding, and flow redistribution.

\subsubsection{Software-Dependency Evidence}

The Express 4.18.2 case instantiated the frozen core on a real topology
of 71 packages and 128 dependency edges, with seven evidence-backed
affected versions. Express root risk increased from 0.5000 to 0.9975
across the transmission sweep and reached rank one by $\tau=0.50$.
Ranking agreement with the $\tau=1$ reference remained high.

This result supports real-topology structural applicability and expected
sensitivity to transmission. It does not establish that the modeled
index is an exploit probability or a predictor of future incidents.

\subsubsection{Electrical-Network Evidence}

The electrical case verified import and baseline preprocessing for 187
bus records, 177 active directed connections, and 269 channel files.
The workflow handled 38 empty archive members without terminating,
reported mapping and missingness explicitly, identified 54 usable
voltage channels, and measured 98.31\% overall timestamp coverage.

The available magnitudes cover normal operation on November 14--15,
2024, while the requested switching events occurred on November 13.
Consequently, the case supports real-topology construction and baseline
data-pipeline verification, but electrical event-response validation
remains pending.

\subsection{Conclusions}

Within the tested scope, the study concludes that:

\begin{enumerate}[label=\arabic*.,leftmargin=*]
  \item GBBRPM v1 provides a deterministic, bounded, and interpretable
  recursive architecture for comparative risk prioritization on DAGs.
  \item Network structure and heterogeneous relation conditions can
  change prioritization beyond what local disturbance alone indicates.
  \item The model is broadly robust to modest input perturbations, but
  exact membership of the highest-priority set becomes less stable as
  uncertainty increases.
  \item The same frozen propagation rule can be instantiated coherently
  in the tested drainage, software, and electrical network structures.
  This supports cross-domain instantiability, not universal validity.
  \item The SWMM disagreement establishes that GBBRPM should not be used
  as a substitute for hydraulic simulation.
  \item The current electrical evidence verifies topology and baseline
  preprocessing only; event-centered validation cannot be claimed until
  temporally aligned measurements are obtained.
\end{enumerate}

\subsection{Limitations}

GBBRPM v1 is limited to directed acyclic graphs and one-pass propagation.
It does not model cycles, feedback, bidirectional influence, explicit
time evolution, or common-cause dependence correction. The bounded risk
state is a comparative index rather than a calibrated probability.
Domain-specific disturbance, susceptibility, and transmission mappings
remain dependent on the quality and defensibility of their input
evidence.

The drainage case uses controlled networks and simulated hydraulic
reference outcomes rather than a calibrated operational field study.
The software case demonstrates structural applicability and sensitivity
rather than observed incident prediction. The electrical case contains
substantial baseline coverage but no measurements aligned with the
requested event windows. The three cases therefore cannot be pooled as
if they offered identical empirical strength.
